\documentclass[12pt]{article}
\usepackage[utf8]{inputenc}
\usepackage[T1]{fontenc}
\usepackage{hyperref}
\title{Sakizaya Language Resources for K-12 Curriculum Development}
\author{Compiled by Claude Sonnet 5.5 (claude-sonnet-5-5) for Andr\'as Kornai}
\date{}
\begin{document}
\maketitle

\section{Introduction}
Sakizaya (ISO 639-3 szy) is an Austronesian language of Taiwan, classed as critically endangered. Most speakers are over 55. It was long treated as a dialect of Amis and was recognised as a separate language in the 2000s. This report is based on a single round of web searches, in English only.

\section{Primary Educational Resources}
The Wikipedia article cites textbooks from the Council of Indigenous Peoples (CIP). I did not see a URL for them. The Alilin e-book collection of the Council (alilin.apc.gov.tw) is linked from Omniglot, but I did not verify that it contains Sakizaya titles.

\section{Linguistic and Reference Materials}
Overviews are at \url{https://en.wikipedia.org/wiki/Sakizaya_language} and \url{https://omniglot.com/writing/sakizaya.htm}. Melissa Shih-hui Lin's 2011 \emph{Sakizaya dictionary}, a report for the CIP dictionary project, is listed at \url{https://glottolog.org/resource/reference/id/695982}. I could not confirm that it is publicly available. A Chinese-language grammar, \emph{Introduction to Sakizaya Grammar}, is mentioned on Omniglot, and I found no publisher details for it.

\section{Religious and Community Resources}
None found.

\section{Digital and Audio Resources}
None confirmed.

\section{Curriculum Development Context}
The language is used by adults only, and no school use is known. Searches in Chinese, which I did not run, would probably turn up more official materials.
\end{document}